\documentclass[12pt]{article}
\usepackage[utf8]{inputenc}
\usepackage[T1]{fontenc}
\usepackage{hyperref}
\title{Wambule Language Resources for K-12 Curriculum Development}
\author{Compiled by Claude Sonnet 5.5 (claude-sonnet-5-5) for Andr\'as Kornai}
\date{}
\begin{document}
\maketitle

\section{Introduction}
Wambule is a Kiranti (Sino-Tibetan) language of the Wambule Rai in eastern Nepal. It is spoken in the Okhaldhunga, Khotang, Udayapur and Sindhuli districts. Speaker counts differ between sources: about 14,000 in the 2011 census according to Wikipedia, and 4,471 in the 2001 census according to another source. Omniglot reports that it is written in Devanagari and taught in primary schools.

\section{Primary Educational Resources}
I found no named Wambule textbooks or primers in search results. Omniglot's report on teaching in primary schools is the only educational evidence: \url{https://omniglot.com/writing/wambule.htm}

\section{Linguistic and Reference Materials}
\begin{itemize}
\item Jean Robert Opgenort, \emph{A Grammar of Wambule} (2004), about 900 pages, with a lexicon and texts: \url{https://glottolog.org/resource/reference/id/120574}. Library record: \url{https://www.oeaw.ac.at/resources/Record/993582211504498}
\item Krishna Prasad Chalise, 2015, sociolinguistic survey of Wambule for the Linguistic Survey of Nepal, Tribhuvan University. I saw no confirmed URL for it.
\item Buddhi Bahadur Rai, 2014 M.Ed. thesis comparing subject-verb agreement in Wambule and English. I saw no confirmed URL for it.
\item Wikipedia overview: \url{https://en.wikipedia.org/wiki/Wambule_language}
\end{itemize}
I found no standalone online Wambule dictionary.

\section{Religious and Community Resources}
I found no scripture or community organisation pages in the search results.

\section{Digital and Audio Resources}
I found no audio, video or corpus resources in the search results.

\section{Curriculum Development Context}
Wambule is stable, with a large speaker base for a minority language, and some schools already teach it as a subject. The Opgenort grammar and lexicon are the main reference material. Curriculum work should find out which textbooks the schools currently use, since the searches did not show them.

\end{document}